\documentclass{resume}
% \usepackage{zh_CN-Adobefonts_external} % Simplified Chinese Support using external fonts (./fonts/zh_CN-Adobe/)
% \usepackage{zh_CN-Adobefonts_internal} % Simplified Chinese Support using system fonts

\newcommand{\SoftwareEngineer}{Software Engineer}
\newcommand{\DataScientist}{Data Scientist}
\newcommand{\ResearchFellow}{Postdoctoral Research Fellow}

\usepackage{xcolor}
\begin{document}
\pagenumbering{gobble} % suppress displaying page number

\name{Guanli Liu}

\basicInfo{
  \email{guanlil1@unimelb.edu.au} \textperiodcentered\
  % \email{liu.guanli22@gmail.com} \textperiodcentered\
  \phone{(+61) 450075953} \textperiodcentered\
  \linkedin[Guanli Liu]{https://www.linkedin.com/in/guanli-liu-11353058/} \textperiodcentered\
  \github[https://github.com/Liuguanli]{https://github.com/Liuguanli}
}

\section{\faUser\ \textbf{Profile}}
I am a postdoctoral researcher specializing in learned spatial indexing and machine learning-driven query optimization. My PhD introduced effective and efficient learned spatial indexing methods (VLDB, ICDE, TKDE). Since finishing in 2023 I have focused on benchmarking RL-based indexing against traditional methods and applying large language models to spatial workloads. Prior to academia, I was a software engineer at Baidu, gaining a strong foundation in algorithms, database systems, and large-scale data processing.

\section{\faFlask\ \textbf{Research Interests}}
\begin{itemize}[parsep=0.4ex]
  \item Learned \& hybrid spatial indexes, cost modeling, and drift-aware benchmarking
  \item GenAI for database tooling (plan summarization, RAG copilots, human-in-the-loop tuning)
  \item Reinforcement learning for query optimization, adaptive storage, and resource allocation
  \item Responsible ML/GenAI evaluation, interpretability, and reproducible experimentation
\end{itemize}

\section{\faGraduationCap\ \textbf{Education}}
\datedsubsection{\textbf{PhD} in Computer Science, \textit{The University of Melbourne}, Australia}{2019 -- 2023}
\datedsubsection{\textbf{M.S.} in Computer Technology, \textit{Northeastern University}, China (GPA: 3.5 / Top 10\%)}{2013 -- 2015}
\datedsubsection{\textbf{B.Eng.} in Software Engineering, \textit{Northeastern University}, China (GPA: 3.4 / Top 20\%)}{2009 -- 2013}

\section{\faUsers\ \textbf{Research and Work Experience}}
\datedsubsection{\textbf{Postdoctoral Research Fellow}, \textbf{The University of Melbourne}, Australia}{Feb. 2024 -- Present}
\begin{itemize}[parsep=0.35ex]
  \item Conduct research under Prof. Renata Borovica-Gajic’s DECRA project on AI for Databases, spanning spatial indexing, data layout optimization, and cardinality estimation.
  \item Designed efficient data layouts, PostgreSQL extensions, and benchmarking frameworks for RL-enhanced indexes; two ICDE 2026 papers under submission.
  \item Built GenAI explainers for optimizer plans (GPT-4 + retrieval) and mentored MS students on reproducibility and benchmarking.
\end{itemize}

\datedsubsection{\textbf{Data Scientist}, \textbf{\href{https://nftdb.ai/}{nftDb \faExternalLink}}, Australia}{Feb. 2023 -- Feb. 2024}
\begin{itemize}[parsep=0.35ex]
  \item Processed NFT transaction data (\textbf{Python}, \textbf{BigQuery}, \textbf{dbt}), managed token wallets, and authored ranking systems using PageRank-like models.
  \item Built core components of the \href{https://databeast.xyz/}{\textbf{databeast}} intelligence platform, detecting wash trading and market anomalies.
  \item Delivered GenAI wallet briefings that combine GPT-4 reasoning with curated graph features, reducing analyst triage time by 40\%.
\end{itemize}

\datedsubsection{\textbf{Research Assistant}, \textbf{The University of Melbourne}, Australia}{Aug. 2022 -- Aug. 2023}
\begin{itemize}[parsep=0.35ex]
  \item \textit{Project One:} Developed an AI-assisted system for reducing reading interruptions using \textbf{GPT API}, \textbf{Google Cloud}, and Midjourney.
  \item \textit{Project Two:} Designed a \textbf{C-language coding style checker} (Python + AST rules) to detect common student errors and auto-suggest fixes.
\end{itemize}

\datedsubsection{\textbf{Software Engineer}, \textbf{Baidu}, China}{Jul. 2015 -- Aug. 2017}
\begin{itemize}[parsep=0.35ex]
  \item Developed features for Baidu’s enterprise IM platform (\href{https://infoflow.baidu.com/}{\textbf{infoflow}}), including new message protocols and voice-assistant modules.
  \item Optimized backend storage/retrieval efficiency by ~20\% and shipped UI enhancements that improved responsiveness.
  \item Partnered with research teams to integrate telemetry and experimentation hooks supporting ML-driven personalization.
\end{itemize}

\section{\faFile\ \textbf{Publications}}
\begin{itemize}[parsep=0.35ex]
  \item \textbf{Guanli Liu}, Renata Borovica-Gajic, Hai Lan, Zhifeng Bao. ``Benchmarking RL-Enhanced Spatial Indices Against Traditional, Advanced, and Learned Counterparts.'' \textit{ICDE}, 2026.
  \item Lankadinee Rathuwadu, \textbf{Guanli Liu}, Christopher Leckie, Renata Borovica-Gajic. ``CoLSE: A Lightweight and Robust Hybrid Learned Model for Single-Table Cardinality Estimation using Joint CDF.'' \textit{ICDE}, 2026.
  \item Kaan Gocmen, \textbf{Guanli Liu}, Renata Borovica-Gajic. ``Advancing Spatial Keyword Queries: From Filters to Unified Vector Embeddings.'' \textit{ADC}, 2025.
  \item Ruiyi Hao, \textbf{Guanli Liu}, Renata Borovica-Gajic. ``LLM-Enhanced Processing of Complex Spatial Queries.'' \textit{ADC}, 2025.
  \item \textbf{Guanli Liu}, Lars Kulik, Christian S. Jensen, Tianyi Li, Renata Borovica-Gajic, Jianzhong Qi. ``Efficient Cost Modeling of Space-filling Curves.'' \textit{VLDB}, 2025.
  \item \textbf{Guanli Liu}, Jianzhong Qi, Lars Kulik, Kazuya Soga, Renata Borovica-Gajic, Benjamin I. P. Rubinstein. ``Efficient Index Learning via Model Reuse and Fine-tuning.'' \textit{ICDEW}, 2023.
  \item \textbf{Guanli Liu}, Jianzhong Qi, Christian S. Jensen, James Bailey, Lars Kulik. ``Efficiently Learning Spatial Indices.'' \textit{ICDE}, 2023.
  \item Jianzhong Qi, \textbf{Guanli Liu}, Christian S. Jensen, Lars Kulik. ``Effectively Learning Spatial Indices.'' \textit{VLDB}, 2020.
  \item Yu Gu, \textbf{Guanli Liu}, Jianzhong Qi, Hongfei Xu, Ge Yu, Rui Zhang. ``The Moving K Diversified Nearest Neighbor Query.'' \textit{TKDE}, 2016.
\end{itemize}

\section{\faGraduationCap\ \textbf{Teaching}}
\begin{itemize}[parsep=0.35ex]
  \item \textbf{INFO20003 – Database Systems (UoM)}: Prepared 2025 S2 teaching materials; will deliver tutorials/lectures this upcoming semester.
  \item \textbf{COMP90018 – Android Application Development (UoM)}: Tutor (2019–2023) covering tutorials, student support, and marking.
  \item \textbf{COMP90041 – Programming and Software Development (UoM)}: Tutor (2019–2023) responsible for tutorials and assessment.
\end{itemize}

\section{\faUsers\ \textbf{Research Service}}
\begin{itemize}[parsep=0.35ex]
  \item \textbf{Conference Reviewer:} CIKM 2024, VLDB 2025 (External), VLDB 2026 (Shadow), KDD 2025 (Excellent Reviewer).
  \item \textbf{Journal Reviewer:} Transactions on Spatial Algorithms and Systems (TSAS) since 2022.
\end{itemize}

\section{\faBriefcase\ \textbf{Supervision}}
\begin{itemize}[parsep=0.35ex]
  \item \textbf{Kaan Gocmen} (Master’s, UoM) — Advancing Spatial Keyword Queries; publication at \textit{ADC 2025}.
  \item \textbf{Ruiyi Hao} (Master’s, UoM) — LLM for Spatial Queries; publication at \textit{ADC 2025}.
\end{itemize}

\section{\faCogs\ \textbf{Engineering Skills}}
\begin{itemize}[parsep=0.35ex]
  \item \textbf{Programming:} Python, Java, C++ with strong design/algorithms background.
  \item \textbf{Data Management:} PostgreSQL, MySQL, MongoDB, BigQuery; PostgreSQL extension development.
  \item \textbf{Cloud \& Tooling:} Google Cloud Platform, Docker/Kubernetes, GitHub Actions, benchmarking harnesses.
  \item \textbf{Machine Learning:} TensorFlow, PyTorch, TorchLib, scikit-learn, RL environments, GenAI APIs.
\end{itemize}

\end{document}
